\documentclass[12pt]{article}
\usepackage{amsmath,amssymb,amsthm}
\usepackage{geometry}
\geometry{margin=1in}
\usepackage{hyperref}
\usepackage{booktabs}
\usepackage{authblk}
\usepackage[T1]{fontenc}

\newtheorem{theorem}{Theorem}
\newtheorem{proposition}[theorem]{Proposition}
\newtheorem{corollary}[theorem]{Corollary}
\theoremstyle{definition}
\newtheorem{definition}[theorem]{Definition}
\newtheorem{remark}[theorem]{Remark}

\title{Charge Conjugation as Orientation Reversal\\
in Hyperbolic Flavor Geometry}

\author[1]{Marvin L. Gentry}
\affil[1]{Independent Researcher, \texttt{drmlgentry@protonmail.com} \\ ORCID: 0009-0006-4550-2663}
\date{\today}

\begin{document}
\maketitle

\begin{abstract}
We identify the charge conjugation operation $\mathcal{C}$ with
orientation reversal of the compact hyperbolic 3-manifolds underlying
the hyperbolic flavor geometry framework. The lepton and quark sector
manifolds $m003$ and $m006$ are both confirmed chiral---neither is
isometric to its mirror image---via two independent methods: the
Chern--Simons invariant, which changes sign under orientation reversal,
and direct verification that \texttt{is\_amphicheiral() = False} in
SnapPy, with symmetry groups $\mathbb{Z}/2\oplus\mathbb{Z}/2$
containing no orientation-reversing elements. Under the identification
$\mathcal{C}: M\mapsto\bar{M}$, the volume and first homology are
preserved (same mass spectrum and generation structure), while the
Chern--Simons invariant changes sign, distinguishing particle from
antiparticle. For the lepton sector: the electron corresponds to
$m003$ with $\mathrm{CS}=\tfrac{1}{4}$ and the positron to
$\overline{m003}$ with $\mathrm{CS}=\tfrac{3}{4}$. The $\mathbb{Z}/5$
torsion structure admits a CP-odd phase $e^{2\pi i/5}$ that breaks the
$M\leftrightarrow\bar{M}$ symmetry, providing a geometric origin for
matter--antimatter asymmetry. All results are computed from census data
using SnapPy~\cite{SnapPy} and are reproducible.
\end{abstract}

\section{Introduction}
\label{sec:intro}

The Standard Model of particle physics treats particles and antiparticles
as distinct objects related by charge conjugation $\mathcal{C}$, but
provides no geometric origin for this operation. In the hyperbolic flavor
geometry framework~\cite{Gentry:CKM,Gentry:PMNS,Gentry:CP}, the flavor
parameters of the Standard Model arise from the arithmetic and geometry
of compact hyperbolic 3-manifolds. The CKM quark mixing matrix emerges
from the holonomy of the manifold $m006 = \texttt{OrientableClosedCensus[43]}$
($H_1 = \mathbb{Z}/5$, $\mathrm{vol} = 2.029$), the PMNS lepton mixing
matrix from the Meyerhoff manifold $m003 = \texttt{OrientableClosedCensus[1]}$
($H_1 = \mathbb{Z}/5$, $\mathrm{vol} = 0.981$), and CP phases from flat
$\mathrm{U}(1)$ connections classified by $H_1(M) = \mathbb{Z}/5$~\cite{Gentry:CP}.

A natural question arises: where do antiparticles sit in this geometric
framework? In quantum field theory, charge conjugation relates a particle
to its antiparticle, reversing all additive quantum numbers (charge,
lepton number, baryon number) while preserving mass and spin. The mass
spectrum is preserved because particles and antiparticles have identical
masses. In a geometric framework where mass is encoded in the volume and
geodesic spectrum of a manifold, the natural geometric operation
preserving these invariants while distinguishing particle from
antiparticle is \emph{orientation reversal}: $M \mapsto \bar{M}$, where
$\bar{M}$ denotes $M$ with its orientation reversed.

This identification is non-trivial only if the manifolds are chiral ---
that is, if $M$ and $\bar{M}$ are geometrically distinct. If $M$ were
amphicheiral (isometric to its mirror image), the operation $M\mapsto
\bar{M}$ would be trivial and could not distinguish particle from
antiparticle. We show that both $m003$ and $m006$ are chiral, confirmed
by two independent methods, making the identification $\mathcal{C}:
M\mapsto\bar{M}$ geometrically meaningful.

The paper is organized as follows. Section~\ref{sec:chirality}
establishes the chirality of the flavor manifolds. Section~\ref{sec:conj}
defines geometric charge conjugation and identifies its transformation
properties. Section~\ref{sec:asymmetry} discusses the connection to
matter--antimatter asymmetry. Section~\ref{sec:predictions} states
falsifiable predictions. Section~\ref{sec:conclusion} concludes.

\section{Chirality of the flavor manifolds}
\label{sec:chirality}

A compact hyperbolic 3-manifold $M$ is \emph{chiral} if it is not
isometric to its orientation-reversed mirror image $\bar{M}$; it is
\emph{amphicheiral} if it is isometric to $\bar{M}$.

\subsection{The Chern--Simons invariant as a chirality detector}

The Chern--Simons invariant $\mathrm{CS}(M) \in \mathbb{R}/\mathbb{Z}$
is a topological invariant of oriented hyperbolic 3-manifolds that
transforms as
\begin{equation}
  \mathrm{CS}(\bar{M}) = -\mathrm{CS}(M) \pmod{1}
  \label{eq:cs_reversal}
\end{equation}
under orientation reversal~\cite{Meyerhoff1987}. If
$\mathrm{CS}(M) \not\equiv 0$ and $\mathrm{CS}(M) \not\equiv
\tfrac{1}{2} \pmod{1}$, then $\mathrm{CS}(M) \neq \mathrm{CS}(\bar{M})$
and $M$ cannot be isometric to $\bar{M}$: the manifold is chiral.

For the lepton sector manifold $m003$:
\begin{equation}
  \mathrm{CS}(m003) = \tfrac{1}{4}, \qquad
  \mathrm{CS}(\overline{m003}) = -\tfrac{1}{4} \equiv \tfrac{3}{4}
  \pmod{1}.
  \label{eq:cs_m003}
\end{equation}
Since $\tfrac{1}{4} \neq \tfrac{3}{4}$, $m003$ and $\overline{m003}$
are geometrically distinct: $m003$ is chiral.

For the quark sector manifold $m006$:
\begin{equation}
  \mathrm{CS}(m006) = -0.114137, \qquad
  \mathrm{CS}(\overline{m006}) = +0.114137.
  \label{eq:cs_m006}
\end{equation}
Since these values differ, $m006$ is chiral.
(The value for $m006$ is given to six decimal places; it is conjectured
to be irrational, consistent with the non-arithmetic status of $m006$.)

\subsection{Direct verification via SnapPy}

Independent confirmation is obtained from the symmetry group of each
manifold. A compact hyperbolic 3-manifold is amphicheiral if and only
if its isometry group contains orientation-reversing elements. Using
SnapPy~\cite{SnapPy}:

\begin{center}
\begin{tabular}{lccc}
\toprule
Manifold & \texttt{is\_amphicheiral()} & Symmetry group & Order \\
\midrule
$m003$ (\texttt{OCC}[1]) & \texttt{False} &
  $\mathbb{Z}/2\oplus\mathbb{Z}/2$ & 4 \\
$m006$ (\texttt{OCC}[43]) & \texttt{False} &
  $\mathbb{Z}/2\oplus\mathbb{Z}/2$ & 4 \\
\bottomrule
\end{tabular}
\end{center}

Both symmetry groups have order 4 and consist entirely of
orientation-preserving isometries. No orientation-reversing element
exists, confirming chirality independently of the Chern--Simons
computation. The two methods agree: both $m003$ and $m006$ are chiral.

\section{Geometric charge conjugation}
\label{sec:conj}

\begin{definition}[Geometric charge conjugation]
\label{def:C}
In the hyperbolic flavor geometry framework, charge conjugation
$\mathcal{C}$ acts as orientation reversal on the underlying manifold:
\[
  \mathcal{C}: M \longmapsto \bar{M}.
\]
\end{definition}

The naturality of this definition rests on the transformation properties
of the geometric invariants that encode physical observables:

\begin{proposition}[Transformation properties of $\mathcal{C}$]
\label{prop:transform}
Under $\mathcal{C}: M \mapsto \bar{M}$:
\begin{align}
  \mathrm{vol}(\bar{M}) &= \mathrm{vol}(M),
    \label{eq:vol_preserved} \\
  H_1(\bar{M}) &\cong H_1(M),
    \label{eq:h1_preserved} \\
  \mathrm{CS}(\bar{M}) &= -\mathrm{CS}(M) \pmod{1}.
    \label{eq:cs_changes}
\end{align}
\end{proposition}

\begin{proof}
Equation~\eqref{eq:vol_preserved}: hyperbolic volume is an orientation-
independent invariant of the Riemannian metric.
Equation~\eqref{eq:h1_preserved}: first homology depends only on the
topological type, not the orientation.
Equation~\eqref{eq:cs_changes}: see~\cite{Meyerhoff1987}.
\end{proof}

\begin{corollary}[Physical interpretation]
\label{cor:physics}
Under $\mathcal{C}$: the mass spectrum is preserved
(vol encodes mass scale), the generation structure is preserved
($H_1 = \mathbb{Z}/5$ encodes three generations), and the
Chern--Simons invariant changes sign, distinguishing particle
from antiparticle.
\end{corollary}

\subsection{Particle--antiparticle assignments}

The identification yields concrete geometric assignments for the
elementary fermions and their antiparticles:

\begin{center}
\begin{tabular}{llcc}
\toprule
Sector & Particle / Antiparticle & Manifold & $\mathrm{CS}$ \\
\midrule
Lepton & Electron $e^-$ & $m003$ & $+\tfrac{1}{4}$ \\
       & Positron $e^+$ & $\overline{m003}$ & $+\tfrac{3}{4}$\footnotemark \\
\midrule
Quark  & Quark $q$      & $m006$ & $-0.1141$ \\
       & Antiquark $\bar{q}$ & $\overline{m006}$ & $+0.1141$ \\
\bottomrule
\end{tabular}
\end{center}
\footnotetext{Chern--Simons invariants are taken modulo~1 in the range
$[0,1)$. Hence $-\tfrac{1}{4} \equiv \tfrac{3}{4}$.}

The rational value $\mathrm{CS}(e^-) = \tfrac{1}{4}$ for the electron
reflects the arithmetic structure of $m003$: arithmetic hyperbolic
3-manifolds have rational Chern--Simons invariants (up to integer
shifts), while non-arithmetic ones generally have irrational
$\mathrm{CS}$~\cite{Chinburg98}. The irrational value
$\mathrm{CS}(q) = -0.1141$ for the quark is consistent with the
conjectured non-arithmetic status of $m006$.

\subsection{Consistency with CPT}

In quantum field theory, CPT is an exact symmetry. In the geometric
framework, the three discrete operations act as:
\begin{itemize}
  \item $\mathcal{C}$: orientation reversal $M \to \bar{M}$,
  \item $\mathcal{P}$: spatial reflection (acts on the embedding
    space, not the internal manifold),
  \item $\mathcal{T}$: time reversal, equivalent to $\mathcal{CP}$
    by the CPT theorem, i.e., orientation reversal combined with
    spatial reflection.
\end{itemize}
The combination $\mathcal{CPT}$ acts as the identity on the manifold
(two orientation reversals), consistent with CPT invariance.

\section{Matter--antimatter asymmetry}
\label{sec:asymmetry}

If $\mathcal{C}$ is orientation reversal, then a universe containing
matter but not antimatter corresponds geometrically to a preferred
orientation of the vacuum manifold. The $\mathbb{Z}/5$ torsion structure
of $m003$ and $m006$ provides the mechanism for this preference.

The flat $\mathrm{U}(1)$ connections on $M$ are classified by
$H^1(M, \mathrm{U}(1)) \cong \mathrm{Hom}(\mathbb{Z}/5, \mathrm{U}(1))$,
giving holonomies $e^{2\pi ik/5}$ for $k = 0, 1, 2, 3, 4$.
The non-trivial character $k=1$ gives a CP-odd phase $e^{2\pi i/5}$.
Under orientation reversal, this phase transforms as
\[
  e^{2\pi i/5} \longmapsto e^{-2\pi i/5},
\]
breaking the $M \leftrightarrow \bar{M}$ symmetry. This is the geometric
origin of CP violation in the lepton sector, and by extension, a
candidate mechanism for leptogenesis~\cite{Fukugita1986}: the CP-odd
phase $e^{2\pi i/5}$ generates a lepton asymmetry that is converted to
a baryon asymmetry via electroweak sphalerons.

The $\mathbb{Z}/5$ torsion yields a Jarlskog invariant
$J = \tfrac{1}{4}\sin(2\pi/5) \approx 0.0309$, which matches the PDG
value $J = 0.033 \pm 0.003$ to within $6\%$ (see companion
paper~\cite{Gentry:CP}). The actual PMNS CP phase $\delta$ depends on
the parameterization; the geometric origin implies $|\sin\delta| \approx
1$, i.e., $\delta \approx 90^\circ$ or $270^\circ$. The PDG best fit
for normal ordering is $\delta \approx 197^\circ$, which is within
$27^\circ$ of $180^\circ+72^\circ=252^\circ$; the precise mapping
requires a full reparameterisation of the geometric matrix.

It is important to note that the present framework provides a geometric
origin for the \emph{existence} of a CP-odd phase but does not
dynamically explain why the universe selected the orientation $M$ rather
than $\bar{M}$. A dynamical mechanism for orientation selection remains
an open problem.

\section{Falsifiable predictions}
\label{sec:predictions}

The geometric identification $\mathcal{C}: M \mapsto \bar{M}$ yields
the following predictions:

\paragraph{Chern--Simons invariants.}
If the Chern--Simons invariant has physical significance (as in
Chern--Simons TQFT approaches to quantum gravity~\cite{Witten1989}),
the electron sector is characterized by $\mathrm{CS} = \tfrac{1}{4}$
and the positron sector by $\mathrm{CS} = \tfrac{3}{4}$. Any future
geometric derivation of electric charge from 3D topology must reproduce
this assignment.

\paragraph{Arithmetic vs non-arithmetic sectors.}
The rational $\mathrm{CS}(m003) = \tfrac{1}{4}$ implies $m003$ is
arithmetic~\cite{Chinburg98}, while the irrational $\mathrm{CS}(m006)
= -0.1141$ implies $m006$ may be non-arithmetic. If $m006$ is proven
arithmetic, its CS invariant would be rational and the current value
would require reinterpretation.

\paragraph{No new species prediction.}
The framework assigns a unique manifold to each fermion generation via
the $\mathbb{Z}/5$ homology structure. Additional fermion generations
would require manifolds with $H_1 \cong \mathbb{Z}/p$ for $p > 5$ or
more complex torsion structures. The absence of a fourth generation
in the closed census with appropriate properties is consistent with
three generations being the complete set.

\paragraph{CP violation magnitude.}
The Jarlskog invariant $J = \tfrac{1}{4}\sin\tfrac{2\pi}{5} \approx
0.031$ is fixed by the $\mathbb{Z}/5$ torsion structure. If future
measurements shift the PDG value $J = 0.033 \pm 0.003$ by more than
$3\sigma$ away from $0.031$, the geometric CP mechanism would be
disfavored.

\section{Conclusion}
\label{sec:conclusion}

We have established that the two compact hyperbolic 3-manifolds
underlying the hyperbolic flavor geometry framework --- $m003$ (lepton
sector) and $m006$ (quark sector) --- are both chiral, confirmed by the
Chern--Simons invariant and by direct SnapPy computation. This chirality
makes the identification of charge conjugation $\mathcal{C}$ with
orientation reversal $M \mapsto \bar{M}$ geometrically non-trivial and
meaningful. Under this identification, the volume and first homology
(mass spectrum and generation structure) are preserved, while the
Chern--Simons invariant changes sign, distinguishing particle from
antiparticle. The electron corresponds to $m003$ with
$\mathrm{CS} = \tfrac{1}{4}$; the positron corresponds to
$\overline{m003}$ with $\mathrm{CS} = \tfrac{3}{4}$.

The $\mathbb{Z}/5$ torsion structure provides a CP-odd phase
$e^{2\pi i/5}$ that breaks the $M \leftrightarrow \bar{M}$ symmetry,
offering a geometric origin for the matter--antimatter asymmetry of
the universe. The Jarlskog invariant $J \approx 0.031$ derived from
this phase matches the PDG value to within $6\%$.

The identification presented here is a natural consequence of the
chirality of the flavor manifolds and the general principle that
orientation reversal is the discrete geometric operation that preserves
metric invariants while reversing orientation-dependent quantities.
Whether this geometric structure reflects a deeper physical principle
or an emergent pattern remains an open question for future work.

\section*{Data Availability}
All computations use SnapPy~\cite{SnapPy} applied to the
\texttt{OrientableClosedCensus}. Scripts are available at
\url{https://github.com/drmlgentry/hyperbolic-flavor-scan}.

\section*{Conflict of Interest}
The author declares no conflicts of interest.

\begin{thebibliography}{99}

\bibitem{PDG2024}
  S.~Navas et al.\ (Particle Data Group),
  Phys.\ Rev.\ D \textbf{110}, 030001 (2024).

\bibitem{Gentry:CKM}
  M.~L.~Gentry,
  ``Quark Mixing from Hyperbolic Geometry: The CKM Matrix from
  Geodesics on a Compact 3-Manifold,''
  SSRN preprint, doi:10.2139/ssrn.6583550 (2026).

\bibitem{Gentry:PMNS}
  M.~L.~Gentry,
  ``Lepton Mixing from the Borel Structure of Hyperbolic Holonomy,''
  SSRN preprint, doi:10.2139/ssrn.6583553 (2026).

\bibitem{Gentry:CP}
  M.~L.~Gentry,
  ``CP Violation from A-Factor Twist Angles of Hyperbolic Holonomy,''
  SSRN preprint, doi:10.2139/ssrn.6583551 (2026).

\bibitem{Meyerhoff1987}
  R.~Meyerhoff,
  ``Density of the Chern--Simons invariant for hyperbolic 3-manifolds,''
  in \emph{Low-dimensional Topology and Kleinian Groups},
  London Math.\ Soc.\ Lecture Note Ser.\ \textbf{112}, 217--239 (1986).

\bibitem{Chinburg98}
  T.~Chinburg, E.~Friedman, K.~N.~Jones, and A.~W.~Reid,
  ``The arithmetic hyperbolic 3-manifold of smallest volume,''
  Ann.\ Scuola Norm.\ Sup.\ Pisa \textbf{30}, 1--40 (2001).

\bibitem{Fukugita1986}
  M.~Fukugita and T.~Yanagida,
  ``Baryogenesis Without Grand Unification,''
  Phys.\ Lett.\ B \textbf{174}, 45 (1986).

\bibitem{Witten1989}
  E.~Witten,
  ``Quantum field theory and the Jones polynomial,''
  Commun.\ Math.\ Phys.\ \textbf{121}, 351--399 (1989).

\bibitem{SnapPy}
  M.~Culler, N.~M.~Dunfield, M.~Goerner, and J.~R.~Weeks,
  ``SnapPy, a computer program for studying the geometry and topology
  of 3-manifolds,'' \url{http://snappy.computop.org}.

\end{thebibliography}

\end{document}